\section{Case Study: Scene-Aware Object Rotation}
\label{sec:case-study}

We return to the original motivating task to validate whether the workflow can produce useful train-ready data and diagnose its own bottlenecks. Scene-aware object rotation is an appropriate first case because it requires stable background context, object identity preservation, viewpoint control, lighting consistency, and grounding consistency. It also exposes an important ambiguity: object rotation is not the same as camera-orbit multiview rendering.

The task-specific rendering contract is:

\contractbox{\texttt{canonical yaw000 state} \(\rightarrow\) \texttt{rotate object yaw} \(\rightarrow\) \texttt{keep scene and camera fixed}}

This contract belongs to the case study rather than the general framework. It prevents two failure modes. First, selecting a separate best render for each target angle can introduce cross-angle drift in material, lighting, scale, or grounding. Second, moving the camera changes the task into camera-orbit multiview rendering. By fixing the scene and camera and rotating only the object yaw, the pair more directly encodes object-level rotation editing.

The baseline dataset uses the canonical front view as the source and non-zero horizontal viewpoints as targets. Each training object contributes seven target views, producing 245 training pairs from 35 training objects. Validation and test partitions are object-disjoint, with 49 validation pairs and 56 test pairs. In the first feedback-guided weak-angle scaling round, referred to internally as R1, evaluation feedback selects weak angles, and the added objects are train-only. The round contributes 60 weak-angle pairs at 90, 180, and 270 degrees, expanding the train set to 305 pairs without contaminating validation or test objects.

This case study therefore connects the two loops. The inner loop constructs and reviews renderable object data. The outer loop uses validation feedback to decide where new data should be added. The result is a concrete train-ready dataset round that can be evaluated without changing the held-out protocol.
